% Part: normal-modal-logic
% Chapter: filtrations
% Section: S5-fmp
%READERNOTE{SOL6-A171: प्रणालीतील प्रत्येक सूत्र प्रतिमानाच्या प्रत्येक जगावर सत्य असणे आणि चौकट-गुणधर्म असणे हे स्वतंत्र आहेत. S5 च्या कोणत्याही प्रतिमानात लक्ष्य सूत्र सत्य असेल तर त्याचा निषेध सिद्ध होत नाही; संपूर्णतेवरून परावर्ती युक्लिडी प्रतिमान मिळते. त्यानंतर वैश्विक प्रतिमान आणि सांत निस्यंदनाचे मूळ पाऊल लागू होते.}

\documentclass[../../../include/open-logic-section]{subfiles}

\begin{document}

\olfileid{nml}{fil}{fmp}

\olsection{\Log{K} आणि \Log{S5} यांना सांत प्रतिमान गुणधर्म आहे}

\begin{defn}
  मोडल तर्कशास्त्राच्या प्रणाली~$\Sigma$ ला
  \emph{सांत प्रतिमान गुणधर्म} आहे असे म्हणतात, जर
  !!a{formula}~$!A$ हे~$\Sigma$ च्या एखाद्या
  प्रतिमानातील एखाद्या जगावर सत्य असेल, तर~$!A$
  हे~$\Sigma$ च्या एखाद्या \emph{सांत} प्रतिमानातील
  एखाद्या जगावरही सत्य असेल. येथे प्रणालीचे प्रतिमान म्हणजे
  त्या प्रणालीतील प्रत्येक सूत्र प्रत्येक जगावर सत्य असलेले प्रतिमान.
\end{defn}

\begin{prop}\ollabel{prop:K-fmp}
  \Log{K} ला सांत प्रतिमान गुणधर्म आहे.
\end{prop}

\begin{proof}
  \Log{K} हा वैध !!{formula}s चा संच आहे;
  म्हणजे कोणतेही प्रतिमान हे~\Log{K} चे प्रतिमान आहे.
  \olref[fil]{thm:filtrations} नुसार,
  $\mSat{M}{!A}[w]$ असेल, तर~$!A$ च्या
  उप-!!{formula}s च्या संच~$\Gamma$ मधून मिळणाऱ्या
  $\mModel{M}$ च्या कोणत्याही निस्यंदनासाठी
  $\mSat{M^*}{!A}[{[w]}]$. कोणत्याही !!{formula}
  ची उप-!!{formula}s सांत संख्येनेच असतात; म्हणून
  $\Gamma$ सांत आहे. \olref[fin]{prop:filt-are-finite}
  नुसार $\card{W^*} \le 2^n$, जेथे $n$ ही~$\Gamma$
  मधील !!{formula}s ची संख्या आहे. शिवाय \Log{K}
  प्रतिमानांवर कोणतेही निर्बंध घालत नसल्याने
  $\mModel{M^*}$ हे \Log{K}-प्रतिमान आहे.
\end{proof}

तर्कशास्त्र~\Log{L} याला निस्यंदनांद्वारे
सांत प्रतिमान गुणधर्म आहे हे दाखवण्यासाठी,
\Log{L}-प्रतिमानाचे निस्यंदन हे स्वतः
\Log{L}-प्रतिमान असणे आवश्यक आहे. यासाठी
अनेकदा बरेच काम करावे लागते आणि कोणतेही
निस्यंदन नेहमीच \Log{L}-प्रतिमान देत नाही.
मात्र वैश्विक प्रतिमानांच्या बाबतीत हे खरे आहे.

\begin{prop}\ollabel{prop:univ-fin}
  $\mClass{U}$ हा वैश्विक प्रतिमानांचा वर्ग
  (\olref[frd][es5]{prop:S5=univ} पाहा) आणि
  $\mClass{U}_\mathrm{Fin}$ हा सर्व सांत वैश्विक
  प्रतिमानांचा वर्ग असू द्या. मग कोणतेही
  !!{formula}~$!A$ हे $\mClass{U}$ मध्ये वैध असेल
  तेव्हा आणि तेव्हाच ते $\mClass{U}_\mathrm{Fin}$
  मध्ये वैध असेल.
\end{prop}

\begin{proof}
  सांत वैश्विक प्रतिमाने ही वैश्विक प्रतिमानेच
  असल्याने डावीकडून उजवीकडे जाणारा निष्कर्ष
  तात्काळ मिळतो. उलट दिशेसाठी, वैश्विक
  प्रतिमान~$\mModel{M}$ मधील एखाद्या जगावर~$w$
  $!A$ असत्य आहे असे समजा. $\Gamma$ मध्ये $!A$
  आणि त्याची सर्व उपसूत्रे घ्या; स्पष्टपणे $\Gamma$
  सांत आहे. $\mModel{M}$ चे एखादे निस्यंदन~$\mModel{M^*}$
  घ्या. \olref[fin]{prop:filt-are-finite} नुसार
  $\mModel{M^*}$ सांत आहे; आणि
  \olref[fil]{thm:filtrations} नुसार~$\mModel{M^*}$
  मधील $[w]$ येथे $!A$ असत्य आहे. $\mModel{M^*}$
  हेही वैश्विक आहे, एवढे आता पाहायचे आहे:
  कोणतीही $u$ आणि $v$ दिली, तर गृहीतकानुसार
  $Ruv$; आणि
  \olref[fil]{defn:filtration}\olref[fil]{defn:filtration-R}
  नुसार $R^*[u][v]$ सुद्धा.
\end{proof}

\begin{cor}\ollabel{cor:S5fmp}
  \Log{S5} ला सांत प्रतिमान गुणधर्म आहे.
\end{cor}

\begin{proof}
  प्रथम, $!A$ हे \Log{S5} च्या कोणत्याही प्रतिमानातील एखाद्या
  जगावर सत्य आहे असे समजा. मग $\lnot!A$ हे \Log{S5} मध्ये
  !!{derivable} नसते; अन्यथा ते त्या प्रतिमानाच्या प्रत्येक जगावर
  सत्य असते. \olref[com][fra]{thm:generaldet} नुसार $\lnot!A$
  हे परावर्ती युक्लिडी प्रतिमानांच्या वर्गात वैध नसते. म्हणून $!A$
  हे त्या वर्गातील एखाद्या प्रतिमानाच्या एखाद्या जगावर सत्य असते.
  \olref[frd][es5]{prop:S5=univ} नुसार,
  $!A$ हे एखाद्या परावर्ती व युक्लिडी प्रतिमानातील
  एखाद्या जगावर सत्य असेल, तर ते एखाद्या वैश्विक
  प्रतिमानातील जगावर सत्य असते.
  \olref{prop:univ-fin} च्या सिद्धतेतील निस्यंदन वापरल्यावर
  ते एखाद्या सांत
  वैश्विक प्रतिमानातील जगावर सत्य असते
  (म्हणजे त्या प्रतिमानाचे~$!A$ च्या उप-!!{formula}s
  च्या संचातून मिळणारे निस्यंदन).
  प्रत्येक वैश्विक प्रतिमान परावर्ती व युक्लिडीही असते;
  म्हणून $!A$ हे सांत परावर्ती युक्लिडी प्रतिमानातील
  एखाद्या जगावर सत्य असते.
\end{proof}

\begin{prob}
  सर्वत्र उत्तरवर्ती किंवा परावर्ती प्रतिमानाचे
  कोणतेही निस्यंदन अनुक्रमे सर्वत्र उत्तरवर्ती किंवा
  परावर्ती असते, हे दाखवा.
\end{prob}

\begin{prob}
  सममित (संक्रमक, युक्लिडी) प्रतिमानाचे
  सममित नसलेले (संक्रमक नसलेले, युक्लिडी नसलेले)
  निस्यंदन शोधा.
\end{prob}

\end{document}
